% GENERATED FILE. Edit canonical lessons, then run npm run generate:books.
% canonical-source-hash: fnv1a64:15d152ff7b9d23bc
% canonical-lessons: JA-W143-my-own-line

\chapter{Writing: My Own Line}
\label{ch:ja-writing-my-own-line}
\hlchaptermodality{143}

\begin{chapteropening}
\textbf{By the end of this chapter:} \emph{I can write one line of Japanese about my day from a closed supply of known words, with nothing to copy, and write every small sign small.}

Complete JA-W143-my-own-line: a time word, a polite verb and a full stop, chosen rather than copied, with the small signs checked before anything else.
\end{chapteropening}

\section[(my own line)]{My own line — the first Japanese with no model}

\label{lesson:JA-W143-my-own-line}



\begin{quote}
Everything you have written so far had something to reach for: a sign to trace, a word to copy, a word someone said.

This does not. You choose the words, from a supply that is closed: only what this book has taught you.
\end{quote}

\subsection*{Writing — controlled composition}
Write \textbf{one line} about your own day, in hiragana. A time word first, then what you do, ending on a polite verb in \textbf{\ja{ます}} and a small circle, \textbf{\ja{。}}, after it. You do not need to say who; you learned in chapter 138 that Japanese leaves that out when it is clear. If you mark a topic with \emph{wa}, write it \textbf{\ja{は}}.

Before you start, one thing to decide on purpose rather than by accident: \textbf{a small sign is a different sign, not a smaller copy.}

\textbf{\ja{っ}} is not \textbf{\ja{つ}} written lazily. It is one held beat, and it sits small and low in its square. Write \textbf{\ja{きつぷ}} for \textbf{\ja{きっぷ}}, a ticket, and you have written a word that does not exist. The same goes for \textbf{\ja{ゅ}} in \textbf{\ja{しゅみ}}: full size, it reads as two beats instead of one.

Two more checks belong to the hand, not the eye. Give every sign \textbf{the same size of square}, small signs included. And write each sign in the \textbf{stroke order} you learned: the order does not show on the finished line, but it is what keeps the shape right when you write fast, which the next lesson will ask you to do.

Write the line. Then read it back and check the \textbf{small signs} before you check anything else.

\begin{tcolorbox}[breakable,colback=teal!4,colframe=teal!35!black,title={Before you move on}]
Is a small \textbf{\ja{っ}} just \textbf{\ja{つ}}, written smaller? (\textbf{No} — it is another sign, one held beat.)

If your line has a topic marked \textbf{wa}, which sign did you write for it? (\textbf{\ja{は}}, as you learned in the first chapters.)

And what did you choose here that the earlier lessons chose for you? (\textbf{The words.})
\end{tcolorbox}
